\uuid{jAtv}
\exo7id{4158}
\titre{exo7 4158}
\auteur{quercia}
\organisation{exo7}
\datecreate{2010-03-11}
\isIndication{false}
\isCorrection{false}
\chapitre{Fonction de plusieurs variables}
\sousChapitre{Dérivée partielle}
\cpge{mpsi-mp2i.fonctions-deux-variables.b, pcsi.fonctions-deux-variables.b, ptsi.fonctions-deux-variables.b, mpsi-mp2i.fonctions-deux-variables.c, pcsi.fonctions-deux-variables.c, ptsi.fonctions-deux-variables.c, bcpst1.fonctions-deux-variables.c, bcpst1.fonctions-deux-variables.d, mp-mpi.calcul-differentiel-optimisation.a, pt.fonctions-deux-trois-variables.g}
\module{Analyse}
\niveau{L2}
\difficulte{}

\contenu{
\question{
Soit $f : {\R^2} \to \R$ de classe $\mathcal{C}^2$.
Démontrer que :
$$f(a+h,b+k) = f(a,b) + \left( h\frac{\partial f}{\partial x} + k\frac{\partial f}{\partial y} \right)(a,b)
  + \frac 12 \left( h^2\frac{\partial^2 f}{\partial x^2} + 2hk\frac{\partial^2 f}{\partial x \partial y} + k^2\frac{\partial^2 f}{\partial y^2} \right)(a,b)
  +  o(h^2+k^2).$$
}
}
